\clearpage
\section{Binding-freedom of every translated lambda-context}
\label{sec:copying-context-binding}
\termdef{copying-context-binding-certificate}
We now prove the binding-free graph obligation in the paper's proof of
Lemma~11.3: every annotated lambda-context translates to a binding-free graph,
for every hole name and every finite hole allowance. The proof uses the
existing Figure~16 construction, including the prescribed distributors,
middle-block exchange and padding. No term or translation clause changes.
Section~\ref{sec:copying-certified-contexts} supplies the sole-hole certificate
and graph-context congruence with an explicit execution hypothesis.
Addressed-plugging correspondence and term-context congruence remain
separate obligations in Appendix~B.1.

The central observation is directional. A serial graph has no path that
crosses from its right component back into its left component. An operation
prefix has no forbidden starting label. A sharing suffix has no hole that
could be the endpoint of a forbidden path. These facts justify the actual
translation graphs without deleting or identifying distributor edges.

\subsection{Adjacency through finite expansion}


For an expanded edge with disjoint source and target lists, the insertion
map is injective on all vertices. Internal vertices carry the component tag;
input and output coincidences are ruled out by the old edge's absence of a
self-loop; equal ports of the same kind are identified using port injectivity.
Consequently adjacency between two edges in that component is preserved and
reflected: pull their shared vertex back through the injective insertion map.

For two distinct components, any adjacency must cross their boundaries.
A shared internal vertex would identify the component tags. The remaining
possibility is an output port of the first component identified with an
input port of the second. Host injectivity then gives an adjacency between
the corresponding old pattern edges. These lemmas are generic finite-graph
results, independent of the copying alphabet or machine.

\lean{Graph/ExpansionAdjacency}{12}{33}{ExpansionAdjacency.insert_injective}


\lean{Graph/ExpansionAdjacency}{34}{48}{ExpansionAdjacency.adjacent_same_iff}


\lean{Graph/ExpansionAdjacency}{49}{67}{ExpansionAdjacency.adjacent_cross}


\subsection{Depth and ordered parallel assembly}
\termdef{copying-binding-padding}


Padding leaves adjacency unchanged under the existing edge bijection.
It also preserves and reflects both the forbidden-starting-label predicate
and the hole-label predicate: primitives retain their constructor, and
boxes remain boxes. Source witnesses are unchanged. The general pullback
lemma therefore proves shallow binding-freedom in both directions.
Induction on depth extends this equivalence to every box body. Together
with focus-freedom invariance it proves binding-freedom invariance under
one-step raising. Equality transport handles depth casts, induction handles
any finite amount of padding, and combining these proves the result for
padding to a requested common allowance. Unsafe paths cannot be concealed
by changing a representation bound.

\lean{Copying/BindingPadding}{11}{15}{BindingPadding.padding_adjacent}


\lean{Copying/BindingPadding}{16}{25}{BindingPadding.padding_bindingLabel}


\lean{Copying/BindingPadding}{26}{45}{BindingPadding.padding_holeLabel}


\lean{Copying/BindingPadding}{46}{69}{BindingPadding.shallow_raise_iff}


\lean{Copying/BindingPadding}{70}{87}{BindingPadding.paths_raise_iff}


\lean{Copying/BindingPadding}{88}{91}{BindingPadding.raise_iff}


\lean{Copying/BindingPadding}{92}{94}{BindingPadding.castDepth_iff}


\lean{Copying/BindingPadding}{95}{100}{BindingPadding.pad_iff}


\lean{Copying/BindingPadding}{101}{104}{BindingPadding.padTo_iff}


In parallel assembly, adjacent edges have the same component tag. Induction
on ancestry therefore keeps every edge in its starting component and
reflects the ancestry to that component. A proposed forbidden path in the
parallel graph yields its source witness and terminal hole in that same
component, contradicting the component's shallow binding-freedom.
At deeper allowances, each inserted label retains its original box body,
so the recursive condition and focus-freedom are inherited as well.

Interface casts preserve these predicates by equality elimination. Hence
the translated pair of subterms is binding-free when both subterms are:
the middle-block exchange only changes the boundary ordering, and the
preceding parallel theorem allows any typed boundary bijection. This proves
the property for the exact variable-before-atom interface of Figure~16.

\lean{Copying/BindingParallel}{14}{26}{BindingParallel.ancestor_parallel}


\lean{Copying/BindingParallel}{27}{40}{BindingParallel.shallow_parallel}


\lean{Copying/BindingParallel}{41}{49}{BindingParallel.paths_parallel}


\lean{Copying/BindingParallel}{50}{56}{BindingParallel.parallel}


\lean{Copying/BindingParallel}{57}{61}{BindingParallel.castInterfaces_iff}


\lean{Copying/BindingParallel}{62}{68}{BindingParallel.pair}


\subsection{Directed serial assembly}
\termdef{copying-serial-direction}


Neither edge of the two-slot serial pattern self-loops: its external and
internal port constructors distinguish its source and target vertices.
A backward adjacency in the expansion would, by the crossing lemma, give
an adjacency from the right pattern edge to the left one. That would identify
an external output with an external input, which is impossible.

Induction on ancestry proves that a path starting in the right component
stays there. If the two endpoints have the same tag, every intermediate tag
must agree too: a left-to-right crossing could not return to a left endpoint.
The within-component adjacency reflection theorem then reconstructs the
whole ancestry in the original component.

\lean{Copying/BindingSerial}{12}{17}{BindingSerial.serial_nonSelfLoop}


\lean{Copying/BindingSerial}{18}{28}{BindingSerial.serial_no_backward}


\lean{Copying/BindingSerial}{29}{40}{BindingSerial.serial_ancestor_monotone}


\lean{Copying/BindingSerial}{41}{64}{BindingSerial.serial_ancestor_same}


Serial shallow closure requires a condition on paths crossing from left
to right. It suffices that either the left graph has no prohibited starting
labels or the right graph has no holes. Under either alternative, a forbidden
path must have both endpoints on the same side. Directionality puts its first
successor on that side as well, and the ancestry theorem reflects the rest
of the path. Pulling back the source witness contradicts that component's
shallow condition.

The recursive version retains the original body condition below every
expanded edge. Combining it with the existing focus-free serial theorem
gives full binding-freedom. In particular, a focus-free, hole-free suffix
can be appended to any binding-free graph. This is the case needed for
the translation's contraction and weakening distributors. Their labels
need not be operations, and their edges are not removed.

\lean{Copying/BindingSerial}{65}{101}{BindingSerial.shallow_serial}


\lean{Copying/BindingSerial}{102}{113}{BindingSerial.paths_serial}


\lean{Copying/BindingSerial}{114}{122}{BindingSerial.serial}


\lean{Copying/BindingSerial}{123}{133}{BindingSerial.serial_suffix}


\subsection{The complete context-translation certificate}


A graph in the translated language is focus-free and hole-free, so it is
binding-free. An operation corolla is such a graph and has no prohibited
starting label; serial closure therefore permits it as a prefix of any
binding-free argument graph. A box has no outer hole, while its body inherits
the binding-free condition independently at the lower depth.

For abstraction, apply the interface cast to the certified body, box it,
prefix the lambda operation and append the original single-use distributors.
For a unary operation, prefix its operation edge and append those distributors.
For application and every binary operation, use the certified ordered pair,
prefix the operation edge, and append the original two-use distributors.
Each suffix is focus-free and hole-free by its existing language theorem.

\lean{Copying/BindingTranslation}{13}{16}{BindingTranslation.language}


\lean{Copying/BindingTranslation}{17}{24}{BindingTranslation.operation_prefix}


\lean{Copying/BindingTranslation}{25}{34}{BindingTranslation.box}


\lean{Copying/BindingTranslation}{35}{43}{BindingTranslation.abstraction}


\lean{Copying/BindingTranslation}{44}{50}{BindingTranslation.unary}


\lean{Copying/BindingTranslation}{51}{57}{BindingTranslation.application}


\lean{Copying/BindingTranslation}{58}{66}{BindingTranslation.binary}


Induction on an annotated lambda-context now proves the stronger statement:
its constructor interpretation preserves binding-freedom for \emph{any}
binding-free graph supplied at the hole. The hole case is the hypothesis.
Abstraction and unary cases use the constructor lemmas. Application and
binary cases combine the recursive certificate with the hole-free translation
of the fixed subterm. Both branches are padded to the original common maximum
allowance, using the proved padding equivalence, before applying the pair
and operation lemmas.

Finally, supply the isolated hole corolla, whose binding-freedom was already
proved. The result is the actual holed context translation, with no
well-behavedness assumption on the context beyond its typed grammar. This
is the first of the two graph properties used in the paper's proof of
Lemma~11.3. The other is correspondence between translated term filling and
addressed graph plugging; it is not asserted by this certificate.

\lean{Copying/BindingTranslation}{67}{88}{BindingTranslation.context_apply}


\lean{Copying/BindingTranslation}{89}{92}{BindingTranslation.context_graph}


A nested example combines abstraction, application and reference creation,
with the enclosing lambda binding a variable in the hole annotation. Its
translation is certified at every allowance and every hole name by the
general theorem.

A negative example places a term contraction immediately before a term hole.
An explicit input vertex supplies a source of the contraction; the shared
internal vertex is both its target and the hole's source. This is a positive
forbidden path, so the graph is not binding-free. The padding equivalence
proves that the same rejection holds at every larger allowance. Thus the
certificate does not accidentally admit the copying configuration that the
binding-free restriction is intended to exclude.

\lean{Copying/BindingTranslationExamples}{10}{12}{BindingTranslationExamples.nestedReferenceContext}


\lean{Copying/BindingTranslationExamples}{13}{18}{BindingTranslationExamples.nestedReferenceContext_bindingFree}


\lean{Copying/BindingTranslationExamples}{19}{22}{BindingTranslationExamples.contractionBeforeHole}


\lean{Copying/BindingTranslationExamples}{23}{32}{BindingTranslationExamples.contractionBeforeHole_not_bindingFree}


\lean{Copying/BindingTranslationExamples}{33}{36}{BindingTranslationExamples.padded_contractionBeforeHole_not_bindingFree}


The full build and all three axiom audits pass. The binding-free certificate
is removed from the unfinished obligations of Appendix~B.1, item~3. The exact sole-hole certificate is subsequently proved in
Section~\ref{sec:copying-certified-contexts}. That item still requires
addressed-plugging correspondence and the term observation/congruence proof. The copying machine and its execution
instance remain part of item~2.
